\section{Колебания}
\subsection{Силовой подход}
%%%%%%%%%%%%%%%%%%%%%%%%%%%%%%%%%%%%%%%%%%%%%%%%%%%%
%%%%%%%%%%%%%%%%%%%%%%%%%%%%%%%%%%%%%%%%%%%%%%%%%%%%
 \z{На сколько отстанут за сутки маятниковые часы, поднятые на высоту Эвереста?}
%
{$\Delta t = \frac{H}{R} t = 2 \text{ мин.}$}
%%%%%%%%%%%%%%%%%%%%%%%%%%%%%%%%%%%%%%%%%%%%%%%%%%%%
\z{Бусинка массы $m$ с зарядом $q$ может двигаться без трения по натянутой нити длины $2L$, на концах которой закреплены заряды $Q$. Найдите частоту малых колебаний системы.}
%
{$T = 2\pi \sqrt{\frac{\pi \varepsilon_0 m L^3}{q Q}}$}
%%%%%%%%%%%%%%%%%%%%%%%%%%%%%%%%%%%%%%%%%%%%%%%%%%%% 
 \z[{\begin{wrapfigure}[3]{r}{0.3\textwidth}%
		\vspace{-10pt}%
		\begin{flushright}%
 			\includegraphics[width=0.3\textwidth]{10.1.3}%
		\end{flushright}%
	\end{wrapfigure}%
	}]%
	{Доска массы $m$ лежит на двух катках, вращающихся с большой скоростью навстречу друг другу. Расстояние между осями катков $l$, коэффициент трения при скольжении доски по катку $\mu$. Найдите частоту продольных колебаний доски.}
%
{$\omega = \sqrt{\dfrac{2 \mu g}{l}}$}
%%%%%%%%%%%%%%%%%%%%%%%%%%%%%%%%%%%%%%%%%%%%%%%%%%%% 
 \z{Найдите частоту колебаний тонкого обруча радиуса $R$, подвешенного на гвозде. Проскальзывания нет; колебания происходят в плоскости обруча.}
%
{$\omega = \sqrt{\dfrac{g}{2 R}}$}
%%%%%%%%%%%%%%%%%%%%%%%%%%%%%%%%%%%%%%%%%%%%%%%%%%%%	
 \z{После загрузки корабля период колебаний его по вертикали увеличится с $7$ до $7.5$ с. Какова масса груза? Сечение по ватерлинии $S = 500 \text{ м}^2$. Характер вовлечения воды в движение считать не изменившимся при загрузке.}
%
{$m \approx 900 \text{ т.}$}